% =============================================================================
% Introduction — workshop variant. Reframed per GPT-5.5 Pro v5 review (sec. C):
% leads with an ARTIFACT CARD and a one-sentence CLAIM HIERARCHY, then
% explicitly disclaims any leaderboard / "first comprehensive" framing.
% Target: <= 0.75 page.
% =============================================================================
\section{Introduction}
\label{sec:intro}

\paragraph{Artifact card.}
This workshop submission is an \textbf{artifact note} accompanying a release
of: \textit{(i)}~one JSON record per Tinker run with configuration, run ID and
per-step mean reward; \textit{(ii)}~per-step Zero-Variance Fraction (ZVF) and
Gradient Utilization (GU) traces for the runs that logged them;
\textit{(iii)}~a per-run manifest (Appendix~\ref{sec:run-manifest}) and a
per-family configuration table (Table~\ref{tab:configs});
\textit{(iv)}~the driver scripts. The release contains no per-group reward
arrays, rollout JSONL, or Tinker model weights (Section~\ref{sec:limitations}). We do
not claim a ``first comprehensive'' GRPO benchmark, a frontier ranking, or a
definitive algorithm comparison.

\paragraph{Claim hierarchy (one sentence).}
The mechanical ZVF-to-centered-reward-contrast identity is theory; the
five-seed Qwen3-8B GSM8K held-out comparison ($+1.3$\,pp, $p{=}0.26$) is a
null result; the pooled ZVF--last-10 correlation ($r{=}{-}0.769$, $N{=}15$)
is a descriptive finding that is task-confounded and changes sign within
GSM8K ($r{=}{+}0.401$, $N{=}19$); everything else (frontier-scale Tinker rows,
tool-use, group-size sweeps) is B-grade when completed at one seed; C denotes
failed, partial or unstarted cells.

\paragraph{Why this is an artifact note rather than a benchmark.}
Reinforcement-learning post-training now drives much of modern language-model
reasoning and alignment~\citep{ouyang2022training,guo2025deepseekr1nature},
but empirical claims are hard to compare:
PPO~\citep{schulman2017proximal}, GRPO~\citep{shao2024deepseekmath}, and
DPO~\citep{rafailov2024direct} are typically reported on different model
families with different framework defaults, different rollout regimes, and
different reward parsers. Rather than declare a winner, we release the
underlying traces and a candid evidence-grade table
(Table~\ref{tab:evidence-grade}) so that downstream readers can re-pool
or stratify as they see fit. The release separates the 5-seed open-stack TRL arm and the 5-seed
Tinker held-out control from the single-seed closed-backend \tinker{} case
studies; the paper describes what those artifacts do and do not support. Discussion prompts for the workshop are
in Section~\ref{sec:conclusion}.
